% =============================================================================
% Overleaf-ready appendix: Moving-body IBM debugging in cerisse
% Drop this file into an appendix/ directory and \input{} from the main text.
% Adjust section level (\section -> \chapter) to match host document.
% =============================================================================

\section{Debugging the moving-body IBM pipeline: a phased case study}
\label{sec:appendix-ibm-moving-v5}

\subsection{Problem statement}

A 2D diamond-wedge airfoil is immersed in a Mach-2 Euler freestream. After a
quiescent $720\,\mu\mathrm{s}$ settle period, a constant pitch rate
$\omega = 2330~\mathrm{rad\,s^{-1}}$ is prescribed about the quarter-chord,
$(x_p,y_p)=(-0.025,\,0)~\mathrm{m}$. The simulation exercises all three
moving-body components simultaneously: the sharp-interface immersed boundary
method (IBM), the AMReX block-structured adaptive mesh refinement (AMR), and
the rigid-body kinematic update.

The observed symptoms before any modification were:
\begin{itemize}
  \item Static IBM (\texttt{ib.move=0}) ran indefinitely without anomaly.
  \item With \texttt{ib.move=1}, the simulation failed shortly after pitch
        onset.
  \item On \texttt{max\_level=0}, dt collapsed at step $\sim$1933
        ($\sim$1\textdegree{} pitch).
  \item On \texttt{max\_level=1} AMR, dt collapsed at step $\sim$2107
        ($\sim$6--9\textdegree{} pitch).
\end{itemize}

\subsection{Debug timeline}

\subsubsection{Phase 1: AMR/\texttt{post\_regrid} hypothesis --- falsified}

An initial hypothesis attributed the failure to inter-level state
contamination via \texttt{FillPatch} during \texttt{post\_regrid}. A cross-level
fresh-cell repair pass was implemented inside
\texttt{CNS::post\_regrid}. It triggered too early, corrupting the otherwise
healthy settle phase, and was reverted.

The decisive experiment was a \texttt{max\_level=0} isolation run. Level-0 still
crashed at step $\sim$1933, proving that AMR was not the sole (nor the
dominant) failure mechanism for the \emph{first} instability.

\subsubsection{Phase 2 (v3): exposed-cell fix via single-donor selection}

The legacy \texttt{fixExposedCells} path averaged the four (2D) or six (3D)
immediate conservative-variable neighbours of a freshly exposed fluid cell.
In supersonic flow across an airfoil leading edge, this blends pre-shock and
post-shock states and routinely produces non-physical (negative-pressure)
initial values.

The fix replaces Pass~1 with a \emph{single-donor} selection using a
directional flow score
\begin{equation}
  \mathcal{S}_j \;=\; -\bigl(u_j\,\delta i + v_j\,\delta j\bigr),
  \label{eq:score}
\end{equation}
where $(\delta i, \delta j)$ is the donor offset from the fresh cell and
$(u_j,v_j)$ is the donor's velocity. Positive $\mathcal{S}_j$ denotes a
donor whose flow advects \emph{into} the fresh cell and is therefore
causally upstream. Ties and negative-score cases fall back to the candidate
whose pressure is closest to the candidate-set mean.

\textbf{Outcome.} On \texttt{max\_level=0}, the simulation ran the full
5000 steps, reaching a cumulative pitch of $\sim$153\textdegree{} with a
stable time-step. Schlieren fields show a coherent detached bow shock and
physically plausible wake shedding. On \texttt{max\_level=1}, the original
step-2107 crash persisted, indicating a second, distinct failure mode.

\subsubsection{Phase 3 (v5): donor validity and near-wall reconstruction}

Three further modifications were applied:
\begin{enumerate}
  \item Pass~1 of \texttt{fixExposedCells} rejects any donor with
        $p \le 0$ or non-finite components.
  \item Passes~2 and 3 of \texttt{fixExposedCells} (ring search and
        nearest-valid-fluid fallback) are also converted to single-donor
        selection with the same score~\eqref{eq:score}.
  \item \texttt{Weno.h}, in the \texttt{eflux\_ibm} path, detects when the
        five-point stencil has been narrowed by a solid/ghost marker
        ($\text{gl}<\text{ng}$ or $\text{gr}<\text{ng}$) and falls back to
        a first-order piecewise-constant Lax--Friedrichs flux at those
        faces.
\end{enumerate}

\textbf{Outcome.} The \texttt{max\_level=1} first-failure at step 2107 was
eliminated. A secondary instability now appears at step $\sim$2397, a net
delay of $\sim$290 advancement steps over the pre-patch baseline. The min
density at the new first-BAD step was $\rho_{\min}\approx 1.61$ (far from
vacuum), confirming that the low-density amplification pathway had been
effectively suppressed.

\subsubsection{Phase 4 (v6): widened band + positivity check --- regression}

A further patch attempted to dilate the IB-aware low-order region and to
add a face-level positivity fallback (reject a reconstruction if the
reconstructed LLF$^{\pm}$ density component is below a floor). The v6 run
\emph{regressed}: the first BAD event moved from step 2397 to step 2325
(72 steps earlier). The minimum density at the failure returned to
$\rho_{\min}\sim 10^{-13}$, matching the low-density amplification pattern
seen in v3.

The likely cause is that the widened low-order band introduced excessive
near-wall dissipation that altered the wake dynamics, triggering a
distinct failure mode earlier. We did not pursue further heuristic
patching; v5 was frozen as the best-verified baseline.

\subsection{Frozen baseline}

\begin{itemize}
  \item \textbf{Repository}: \texttt{melampous/cerisse}.
  \item \textbf{Branch}: \texttt{baseline/ibm-v5-clean}.
  \item \textbf{Tag}: \texttt{v5-frozen-debug} (commit
        \texttt{dc70f8a4}) --- preserves the full debug-instrumented state.
\end{itemize}

\textbf{Scope note.} The \texttt{near\_ib} first-order fallback in
\texttt{Weno.h} is active on \emph{all} IBM cases on this branch, not only
moving-body cases. The branch is therefore intentionally dedicated to
moving-FSI debugging, \emph{not} a general-purpose IBM baseline. Merging
into \texttt{ibm\_fsi} should be gated behind either a static-IBM
regression sweep or a runtime switch.

\subsection{Current capability and limitations}

\begin{table}[h]
\centering
\begin{tabular}{lll}
\toprule
Configuration & Status & Note \\
\midrule
Static IBM (\texttt{ib.move=0})              & OK          & Unchanged \\
Moving, \texttt{max\_level=0}                & OK          & 5000 steps, $\sim$153\textdegree{} \\
Moving, \texttt{max\_level=1} AMR            & Partial     & First-failure fixed; long-time not fully reliable \\
Moving, \texttt{max\_level}$\ge 2$           & Untested    & --- \\
Moving, viscous (Navier--Stokes)             & Untested    & --- \\
\bottomrule
\end{tabular}
\caption{Verified capabilities of the \texttt{baseline/ibm-v5-clean} branch.}
\end{table}

\textbf{On the physical status of the current solution.}
\begin{description}
  \item[Level 0:] Physically acceptable and engineering-usable. Suitable for
    visualisation, trend analysis, and method-development baselines. It is
    \emph{not} a fully converged or publication-verified result.
  \item[Level 1:] Substantially improved over the pre-patch baseline, but
    long-time AMR solutions remain not fully reliable. The result
    demonstrates that the first-failure donor-contamination mechanism has
    been eliminated and provides direction for subsequent method work. It
    is \emph{not} suitable as a final physical solution, quantitative
    conclusion, or production result.
\end{description}

\subsection{Secondary root cause: localised to module, not to face}

The pipeline tripwires (stages T1 through T4; see
\texttt{src/ibm/ibm\_tripwire.h}) showed that, at the first BAD step
$\sim$2373 on level=1, the values reported at T1 (step entry), T2 (post
\texttt{rebuildIBM}), and T3 (post \texttt{fixExposedCells}) are
\emph{numerically identical}. This implies that the bad state was already
present at the start of the step and therefore born during the
\emph{previous} step's advance, i.e.\ in the
\texttt{compute\_rhs}/reconstruction/update path rather than in
\texttt{rebuildIBM} or \texttt{fixExposedCells}.

\textbf{Evidence bound.} The present evidence suffices to localise the
secondary instability to the \texttt{compute\_rhs} path, but does \emph{not}
suffice to establish a single formula, a single face, or a single
coarse-fine exchange as the sole responsible sub-mechanism.

The observed behaviour is consistent with the following qualitative
picture:
\begin{itemize}
  \item At larger angles of attack the suction side of the airfoil develops
        a strong expansion region.
  \item High-order WENO-Z5 reconstruction in the vicinity of the IBM ghost
        layer or at a coarse-fine interface can overshoot there.
  \item An overshoot produces a negative-pressure or ill-conditioned
        high-momentum cell.
  \item Subsequent AMR synchronisation propagates the bad state across
        levels.
\end{itemize}

\subsection{Next-phase technical directions (prioritised)}

\begin{enumerate}
  \item \textbf{Primitive-space MUSCL with limiter in the IB neighbourhood.}
    Retain WENO-Z5 away from the wall. In the IB-adjacent band, replace
    characteristic WENO with a primitive-space MUSCL reconstruction using a
    TVD limiter (e.g.\ minmod or MC). This removes the ill-conditioned
    characteristic / Roe transform in a region where the assumptions behind
    it are weak.
  \item \textbf{Face-local positivity-preserving fallback.}
    Not a post-RK clip of cell averages. Check the reconstructed left/right
    primitive states at each face; if $\rho$ or $p$ is unphysical, revert
    that face to a plain Lax--Friedrichs or Rusanov flux.
  \item \textbf{Systematic IBM/ghost consistency.} Primitive-variable
    reconstruction in the ghost-adjacent stencil; explicit handling of
    moving-wall velocity in the ghost-point image; a coarse-fine exchange
    that is aware of the IB band.
  \item \textbf{Backup options (not recommended as first attempt):}
    Riemann-based ghost-fluid method after Sambasivan--Udaykumar (2009);
    geometry-rebuild scheduling by accumulated rotation threshold rather
    than every step; cut-cell IBM (Seo--Mittal, 2011), which would amount to
    a full rewrite of the IBM module and is not presently justified.
\end{enumerate}

\subsection{Engineering lessons}

\begin{itemize}
  \item The single most productive experiment was the \texttt{max\_level=0}
        isolation run, which removed AMR from the suspect list early.
  \item Pipeline tripwires (a handful of inline min/max scans at
        well-chosen stages) outperformed log inspection for locating the
        origin of a corrupted state.
  \item Heuristic patches showed rapidly diminishing returns: v3 to v5
        bought $\sim$290 steps; v5 to v6 regressed by 72 steps. Knowing when
        to stop was as important as the fixes themselves.
  \item Every version, including regressions, was kept in the debug record.
        Without this, later readers could have mistaken the regressed v6 for
        an improvement over v5.
  \item External peer review --- even asynchronous and not tool-integrated
        --- corrected multiple conceptual errors (notably, conflating a
        ghost-cell wall trace with a freshly-opened cell's volume average)
        and was decisive in avoiding over-aggressive patching.
\end{itemize}

\subsection{Repository artefacts}

\begin{itemize}
  \item \texttt{src/ibm/ibm\_solver.h} --- \texttt{fixExposedCells} rewrite.
  \item \texttt{src/rhs/Weno.h} --- \texttt{eflux\_ibm} near\_ib fallback.
  \item \texttt{src/ibm/ibm\_tripwire.h} --- four-stage pipeline scanner,
        disabled by default.
  \item \texttt{exm/fsi/airfoil\_pitchup\_euler/} --- test case.
  \item \texttt{docs/debug\_moving\_ibm\_v5\_summary.md} --- Chinese-language
        debug narrative with commit/tag references and test-data paths.
  \item \texttt{docs/BRANCH\_NOTES\_baseline-ibm-v5-clean.md} --- branch
        README.
\end{itemize}
